\begingroup
\renewcommand{\arraystretch}{1.45}
\fancyhead[L]{\small LEMNISCATIC FUNCTIONS II.}
\Needspace{20\baselineskip}
\CurvaArticle{17}{[17.]}
\[
\begin{aligned}
&1+\Bigl(\frac12\Bigr)^3x^4+\Bigl(\frac12.\frac34\Bigr)^3x^8+
\\[7pt]&\quad{}\Bigl(\frac12.\frac34.\frac56\Bigr)^3x^{12}+\ .\ .\ .\\[5pt]
&=\Bigl\{1+\Bigl(\frac14\Bigr)^2x^4+\Bigl(\frac14.\frac58\Bigr)^2x^8+
\\[7pt]&\quad{}\Bigl(\frac14.\frac58.\frac9{12}\Bigr)^2x^{12}+\ .\ .\Bigr\}^2
\end{aligned}
\]
\Needspace{12\baselineskip}
\textit{Proof.} Set
\[
1+\left(\frac14\right)^2x^4+\left(\frac14.\frac58\right)^2x^8+\ .\ .\ .=t
\]
\Needspace{12\baselineskip}
Then, writing
\[
x(x^4-1)\frac{\mathrm{dd}t}{\mathrm dx^2}-(3x^4-1)\frac{\mathrm dt}{\mathrm dx}+x^3t=R,\quad R=0
\]
\Needspace{16\baselineskip}
it also follows that
\[
\begin{aligned}&0=x\frac{\mathrm dR}{\mathrm dx}+R;\\[7pt]&\text{and also}\quad
0=2xt\frac{\mathrm dR}{\mathrm dx}+2\left(t+3x\frac{\mathrm dt}{\mathrm dx}\right)R\end{aligned}
\]
\Needspace{24\baselineskip}
Expanding then gives
\[
\begin{aligned}
0={}&xx(x^4-1)\left(2t\frac{\mathrm d^3t}{\mathrm dx^3}+6\frac{\mathrm dt}{\mathrm dx}\frac{\mathrm{dd}t}{\mathrm dx^2}\right)
\\[7pt]&\quad{}+3x(3x^4-1)\left(2t\frac{\mathrm{dd}t}{\mathrm dx^2}+2\frac{\mathrm dt}{\mathrm dx}\frac{\mathrm dt}{\mathrm dx}\right)\\[5pt]
&+(19x^4-1)2t\frac{\mathrm dt}{\mathrm dx}+8x^3tt
\end{aligned}
\]
\Needspace{16\baselineskip}
or, on setting $tt=u$
\[
\begin{aligned}0&=xx(x^4-1)\frac{\mathrm d^3u}{\mathrm dx^3}+3x(3x^4-1)\frac{\mathrm{dd}u}{\mathrm dx^2}
\\[7pt]&\quad{}+(19x^4-1)\frac{\mathrm du}{\mathrm dx}+8x^3u\end{aligned}
\]
\Needspace{12\baselineskip}
As a solution of this equation one finds
\[
u=1+\left(\frac12\right)^3x^4+\left(\frac{1.3}{2.4}\right)^3x^8+\ .\ .
\]

\Needspace{18\baselineskip}
\SourceLeaf{434}

Now, since
\[
t=\frac{\surd 8}{\pi}\int\frac{\mathrm dz}{{}^{4}\!\surd{(1-z^4)(1-x^4z^4)}}
\]
from $z=0$ to $z=1$, one obtains, for $x=1$,
\[
t=\frac{\surd 8}{\pi}.\frac{\varpi}{2}=\frac{\varpi}{\pi}\surd 2
\]
\Needspace{14\baselineskip}
and therefore
\[
1+\left(\frac12\right)^3+\left(\frac{1.3}{2.4}\right)^3+
\left(\frac{1.3.5}{2.4.6}\right)^3+\ .\ .\ .=2\frac{\varpi\varpi}{\pi\pi}
\]
\Needspace{24\baselineskip}
\begin{center}\textit{The same in another way}\end{center}
The value of the series
\[
1+\left(\frac12\right)^3+\left(\frac{1.3}{2.4}\right)^3+
\left(\frac{1.3.5}{2.4.6}\right)^3+\ .\ .
\]
is
\[
=\frac4{\pi\pi}\iint\frac{\mathrm d\varphi\,\mathrm d\psi}{\surd {(1-\cos\varphi^2\cos\psi^2)}}
\]
from $\varphi=0$ to $\varphi=90^0$ and from $\psi=0$ to $\psi=90^0$.
\Needspace{18\baselineskip}
Thus, setting $\cos\varphi\cos\psi=\cos\upsilon$, the same value becomes
\[
=\frac4{\pi\pi}\iint\frac{\mathrm d\varphi\,\mathrm d\upsilon}{\surd {(\cos\varphi^2-\cos\upsilon^2)}}
\]
namely, from $\upsilon=0$ to $\upsilon=90^0$ and from $\varphi=0$ to $\varphi=\upsilon$.
\Needspace{18\baselineskip}
Finally, setting $\varphi+\upsilon=f$, $\upsilon-\varphi=g$, our expression becomes
\[
=\frac2{\pi\pi}\iint\frac{\mathrm df.\mathrm dg}{\surd {\sin f\sin g}}
\]
from $g=0$ to $g=90^0$ and from $f=g$ to $f=180^0-g$.
\Needspace{16\baselineskip}
But it is not difficult to prove that the integral has the same value if taken
from $f=0$ to $f=90^0$ and from $g=0$ to $g=90^0$, from which its value follows
\[
=2\frac{\varpi\varpi}{\pi\pi}
\]
as above.

\clearpage
\endgroup
